{\color{ink}\bfseries 例 5-2 \(\cdot \) 合订本144 \(\cdot \) 图1}\par \textbf{单缓冲：前3块（\(\mu \)s）}\par\vspace{2bp}\begin{tikzpicture}[x=1bp,y=-1bp,>=Stealth,font=\fontsize{8.5}{10.7}\selectfont,every node/.style={inner sep=1bp},line width=.45bp]
\node[anchor=north west,align=left] at (0,12){设备输入};
\draw[fill=blue] (5.0,24) rectangle (30.0,38);
\node[anchor=center,inner sep=0pt] at (17.5,31){1};
\draw[fill=blue] (42.5,24) rectangle (67.5,38);
\node[anchor=center,inner sep=0pt] at (55.0,31){2};
\draw[fill=blue] (80.0,24) rectangle (105.0,38);
\node[anchor=center,inner sep=0pt] at (92.5,31){3};
\node[anchor=north west,align=left] at (0,36){CPU复制、分析};
\draw[fill=orange] (30.0,48) rectangle (42.5,62);
\node[anchor=center,inner sep=0pt] at (36.25,55){1};
\draw[fill=green] (42.5,48) rectangle (55.0,62);
\node[anchor=center,inner sep=0pt] at (48.75,55){1};
\draw[fill=orange] (67.5,48) rectangle (80.0,62);
\node[anchor=center,inner sep=0pt] at (73.75,55){2};
\draw[fill=green] (80.0,48) rectangle (92.5,62);
\node[anchor=center,inner sep=0pt] at (86.25,55){2};
\draw[fill=orange] (105.0,48) rectangle (117.5,62);
\node[anchor=center,inner sep=0pt] at (111.25,55){3};
\draw[fill=green] (117.5,48) rectangle (130.0,62);
\node[anchor=center,inner sep=0pt] at (123.75,55){3};
\node[anchor=north,inner sep=0pt] at (5.0,68){0};
\node[anchor=north,inner sep=0pt] at (30.0,68){100};
\node[anchor=north,inner sep=0pt] at (55.0,68){200};
\node[anchor=north,inner sep=0pt] at (80.0,68){300};
\node[anchor=north,inner sep=0pt] at (105.0,68){400};
\node[anchor=north,inner sep=0pt] at (130.0,68){500};
\end{tikzpicture}\par 蓝\(=\allowbreak{}\)输入，黄\(=\allowbreak{}\)复制，绿\(=\allowbreak{}\)分析；数字为块号。\par 前3块分析完成：500 \(\mu \)s。
